\documentclass[10pt,letterpaper]{article}

\usepackage[T1]{fontenc}
\usepackage[utf8]{inputenc}
\usepackage{lmodern}
\usepackage[margin=0.82in,headheight=14pt]{geometry}
\usepackage{microtype,amsmath,amssymb,array,tabularx,longtable,booktabs,multicol}
\usepackage{graphicx,xcolor,enumitem,fancyhdr,hyperref,titlesec,needspace}
\definecolor{Navy}{HTML}{122A42}
\definecolor{Ocean}{HTML}{147D91}
\definecolor{Soft}{HTML}{E9F2F5}
\definecolor{Ink}{HTML}{263745}
\definecolor{Gray}{HTML}{5F6873}
\hypersetup{colorlinks=true,linkcolor=Ocean,urlcolor=Ocean,bookmarksnumbered=true}
\setlength{\parindent}{0pt}
\emergencystretch=2em
\setlength{\parskip}{5.3pt}
\setlist{topsep=3pt,itemsep=2pt,parsep=0pt,leftmargin=*}
\titleformat{\section}{\large\bfseries\color{Navy}}{\thesection}{0.65em}{}
\titleformat{\subsection}{\normalsize\bfseries\color{Ocean}}{\thesubsection}{0.6em}{}
\titleformat{\subsubsection}{\normalsize\bfseries\color{Navy}}{\thesubsubsection}{0.5em}{}
\pagestyle{fancy}\fancyhf{}
\fancyhead[L]{\textsf{Architecture Documentation \& Engineering}}
\fancyhead[R]{\textsf{\DocCategory}}
\fancyfoot[L]{\footnotesize\textcolor{Gray}{Controlled learning / engineering artifact}}
\fancyfoot[R]{\footnotesize\thepage}
\renewcommand{\headrulewidth}{0.35pt}
\newcommand{\DocCategory}{Suite}
\newcommand{\ArchTitle}[4]{%
\renewcommand{\DocCategory}{#3}%
{\small\sffamily\bfseries\color{Ocean} #3 \hfill #4}\par\vspace{7pt}
{\Large\bfseries\color{Navy}#1}\par\vspace{4pt}
{\normalsize\color{Gray}#2}\par\vspace{11pt}
\textcolor{Ocean}{\rule{\linewidth}{1.15pt}}\vspace{4pt}
}
\newcommand{\Note}[1]{\par\smallskip\noindent\fcolorbox{Ocean}{Soft}{\parbox{0.94\linewidth}{\small\textbf{Engineering note.}~#1}}\par\smallskip}
\newcommand{\Outcome}[1]{\textbf{Outcome.} #1\par}
\newcommand{\Evidence}[1]{\textbf{Required evidence.} #1\par}
\newcommand{\BlankLines}[1]{\par\begingroup\setlength{\parskip}{9pt}\count255=0\loop\ifnum\count255<#1\noindent\rule{\linewidth}{0.35pt}\par\advance\count255 by 1\repeat\endgroup\vspace{4pt}}
\newcommand{\BlankBox}[1]{\noindent\fbox{\parbox[t][#1][t]{0.98\linewidth}{\mbox{}}}\par\smallskip}
\newcommand{\Field}[2]{\Needspace{3\baselineskip}\textbf{#1}\par\textit{\color{Gray}#2}\par\BlankLines{2}}
\newcommand{\FieldBlank}[1]{\Needspace{3\baselineskip}\textbf{#1}\par\BlankLines{2}}
\newcommand{\DocProvenance}{\textit{Source foundation: Clements et al., \emph{Documenting Software Architectures: Views and Beyond}, 2nd ed. (2011). Interpretive examples, exercises, checklists, and templates are original adaptations, not reproductions.}\par}
\newcommand{\SmallHeader}[1]{\smallskip\textbf{\color{Navy}#1}\par}
\newcolumntype{Y}{>{\raggedright\arraybackslash}X}
\newcommand{\Tick}{\(\square\)\,}

\begin{document}
\ArchTitle{Viewpoint Specification: Module Uses and Layering}{Reusable modeling rules, concerns, correspondence requirements, and verification}{Viewpoint Specification}{ARCH-VP-002}
\DocProvenance
\section{Viewpoint purpose and applicability}
This is a reusable viewpoint specification, not a diagram or description of any specific system. Primary concerns: dependency direction, permitted dependencies and cycles. Principal stakeholders: maintainers, modularity reviewers, build engineers.

\section{Model kinds and relation vocabulary}
module nodes, uses directed edges, optionally layered assignment

\section{Conventions and notation}
directed dependency graph and dependency matrix; assign identifiers to all significant elements and disclose legend, scope, environment and exclusions.

\section{Correspondence rules}
\begin{itemize}
\item each forbidden uses edge is visible and requires reviewed exception
\item Cross-view references use stable IDs and never infer equality from similar names.
\item Stakeholder concerns and quality scenarios must link to this viewpoint when used.
\end{itemize}
\section{Analysis and verification}
\begin{itemize}
\item Validate complete definition of relation semantics and correct element typing.
\item Use reviewer questions from the scorecards appropriate to this viewpoint.
\item Compare every critical claim to source-controlled artifacts and/or runtime evidence; mark confidence and timestamps.
\end{itemize}
\section{Known limitations and failure modes}
confusing uses with membership; hidden cycles; unjustified layer crossing

\section{Required view deliverables}
\small
\begin{tabularx}{\linewidth}{@{}YY@{}}
\toprule
\textbf{Item} & \textbf{Acceptance evidence} \\
\midrule
Primary presentation & directed dependency graph and dependency matrix \\
Element and relation catalog & module nodes, uses directed edges, optionally layered assignment \\
Context and variability & boundary, neighbors, configuration choices \\
Rationale and references & ADR IDs, design constraints, requirements and proof links \\
\bottomrule
\end{tabularx}
\normalsize

\end{document}
